\section{Variational setting, dictionaries, and algorithms}
\label{sec:setting}

Let \(X\) be a real reflexive Banach space. The abstract descent results apply to a proper convex
functional \(\E:X\to(-\infty,+\infty]\), with effective domain
\(\operatorname{dom}\E=\{v\in X:\E(v)<+\infty\}\).
Differentiability is assumed on an open neighborhood of the relevant iterates in the specified
energy space, on which \(\E\) is finite and continuously Fr\'echet differentiable. For an energy
with a restricted effective domain, the additional topology and the neighborhood are stated
explicitly. When an argument invokes a Hessian, its stated neighborhood is in addition contained in
the twice differentiable part of the effective domain. Suppose that \(\E\) has a unique minimizer
\(u^*\) in the admissible space. The duality pairing between \(X^*\) and \(X\) is denoted by
\(\ip{r}{v}\), and the energy defect is
\[
  \delta(u)=\E(u)-\E(u^*), \qquad \delta_m=\delta(G_m).
\]
The examples in \Cref{sec:models} satisfy this setting after the Neumann nullspace has been
removed where necessary.

\subsection{Ridge dictionaries and admissibility}

For an integer activation order \(k\geq1\), a raw ridge atom has the form
\[
  \varphi_{w,b}(x)=\relu(w\cdot x+b)^k,
  \qquad (w,b)\in\R^d\times\R.
\]
An admissible dictionary consists of nonzero ridge functions satisfying the constraints of \(X\),
after the corresponding linear projection. We use a symmetric normalized dictionary
\(\D_X\subset X\):
\[
g\in\D_X\Longrightarrow-g\in\D_X,\qquad \norm g_X=1\quad(g\in\D_X).
\]
When density is needed, we also require \(\overline{\spn\D_X}^{\,X}=X\).
We abbreviate \(\D_X\) to \(\D\) unless a different normalization is specified.

For pure-gradient Neumann energies the admissible space is
\[
  X_{p,\circ}=\left\{v\in W^{1,p}(\Omega):\int_\Omega v\dd x=0\right\}.
\]
Each ridge is projected onto $X_{p,\circ}$ before the normalized dictionary is formed.
Since an almost constant ridge may have arbitrarily small projected norm, normalization
can amplify its derivatives. Hence uniform derivative bounds for the normalized
dictionary require a positive lower bound on the projected norms. Zero projected
atoms are excluded from the analytical dictionary; the corresponding nondegeneracy
assumption is stated whenever it is needed.

\begin{definition}[Atomic hull and atomic norm]
\label{def:atomic}
For a symmetric normalized dictionary \(\D\subset X\), define
\[
 K_1(\D)=\overline{\left\{\sum_{j=1}^Jc_jg_j:
 J<\infty,\ g_j\in\D,\ \sum_{j=1}^J\abs{c_j}\leq1\right\}}^{\,X},
\]
and
\[
 \norm{v}_{\mathcal A(\D)}=
 \inf\{B>0:v/B\in K_1(\D)\}.
\]
\end{definition}

For a bounded set \(K\subset X\), its \(n\)-th dyadic entropy number is
\[
 \varepsilon_n(K;X)=\inf\left\{h>0:\exists x_1,\ldots,x_{2^{n-1}}\in X,
 K\subset\bigcup_{j=1}^{2^{n-1}}(x_j+hB_X)\right\},\qquad n\geq1,
\]
where $B_X=\{x\in X:\norm x_X\leq1\}$. Thus the covering balls may have centers in $X$.
This $2^{n-1}$-ball convention is used throughout the manuscript. We write
$A\lesssim B$ when $A\leq CB$ with a constant independent of the iteration index; a
subscript on $\lesssim$ lists the permitted parameter dependence, and $A\simeq B$ denotes the two
corresponding inequalities.

\begin{assumption}[Atomic source]
\label[assumption]{ass:atomic}
There is a finite \(B>0\), with \(B\geq\norm{u^*}_{\mathcal A(\D)}\), such that
\(u^*/B\in K_1(\D)\).
\end{assumption}

This is a property of the target relative to the \emph{theoretically normalized} dictionary. It is not
interchangeable with the coefficient \(\ell^1\)-norm produced by a discretely normalized algorithm. In one
space dimension, \Cref{prop:1dsource} gives a sufficient condition for finite atomic norm.

\begin{proposition}[A one-dimensional \(\relu^k\) source class]
\label{prop:1dsource}
Let \(I=(0,1)\), \(k\geq1\), \(1\leq p<\infty\), and \(X=W^{1,p}(I)\). Suppose that the continuous
dictionary contains the normalized nonzero truncated powers \((x-t)_+^k/\|(\cdot-t)_+^k\|_X\),
\(0\leq t<1\), together with normalized fully active ridge functions spanning \(\mathbb P_k\). If
\(v\in W^{k+1,1}(I)\), then \(v\) belongs to the atomic space and
\[
\norm v_{\mathcal A(\D)}\leq C_{k,p}\left(
 \sum_{j=0}^k\abs{v^{(j)}(0)}+\norm{v^{(k+1)}}_{L^1(I)}\right),
\]
where \(C_{k,p}\) also accounts for the fixed exponent \(p\), the interval, and the normalization
of the raw atoms. It is a continuous-dictionary source constant and is not the coefficient
\(\ell^1\)-norm returned by a discretely normalized numerical solve.
\end{proposition}

The proof is supplied in the supplementary theory file and uses the integral remainder formula. Multivariate source and
entropy inputs are instead supplied by shallow-network variation-space and width theorems
\citep{SiegelXuVariation2023,SiegelXuSharp2024,MaoSiegelXu2026}; their domain and norm hypotheses
must be checked separately for each application.

\begin{lemma}[Projection and renormalization]
\label{lem:projection}
Let \(P:X\to X_0\) be bounded and linear. If a subdictionary \(\D_0\subset\D\) satisfies
\[
 0<c_0\leq\norm{Pg}_{X_0}\leq C_0<\infty\qquad(g\in\D_0),
\]
and \(\widetilde\D_0=\{Pg/\norm{Pg}_{X_0}:g\in\D_0\}\), then
\[
 P K_1(\D_0)\subset C_0K_1(\widetilde\D_0).
\]
The same computation also shows that a bounded linear map increases every entropy radius by at most
\(\norm P\), while the inverse normalization loses at most the factor \(c_0^{-1}\):
\[
 K_1(\widetilde\D_0)\subset c_0^{-1}P K_1(\D_0).
\]
\end{lemma}

\begin{proof}
For an absolute convex combination, write
\(Pg=\norm{Pg}_{X_0}\widetilde g\) term by term and bound the new coefficient sum by \(C_0\).
Pass to the closure for the first claim. Applying \(P\) to an \(h\)-cover gives a
\(\norm P h\)-cover. Conversely,
\(Pg/\norm{Pg}_{X_0}=c_0^{-1}P[(c_0/\norm{Pg}_{X_0})g]\), and the coefficient
\(c_0/\norm{Pg}_{X_0}\) is at most one. Taking absolute convex combinations and closures proves the
reverse inclusion. This statement concerns convex hulls; it does not assert that pointwise
renormalization is Lipschitz near null projected atoms. The image of \(K_1(\D_0)\) is closed because
this bounded closed convex subset of the reflexive space \(X\) is weakly compact, and its image
under \(P\) is weakly compact in \(X_0\).
\end{proof}

\begin{corollary}[Transfer of source and entropy inputs]
\label{cor:projected-source-entropy}
In the setting of \Cref{lem:projection}, suppose \(u^*=Pu^\dagger\),
\(u^\dagger/B\in K_1(\D_0)\), and
\[
 \varepsilon_n(K_1(\D_0);X)\leq C_{\rm ent}n^{-\beta}.
\]
Then
\[
 \frac{u^*}{BC_0}\in K_1(\widetilde\D_0),
 \qquad
 \varepsilon_n(K_1(\widetilde\D_0);X_0)
 \leq c_0^{-1}\norm P\,C_{\rm ent}n^{-\beta}.
\]
Thus a ReLU source/convex-hull entropy theorem on an extended parameter domain passes to the projected,
renormalized feasible dictionary whenever the nondegeneracy bounds in \Cref{lem:projection} hold.
\end{corollary}

\begin{proof}
The first assertion follows from
\(PK_1(\D_0)\subset C_0K_1(\widetilde\D_0)\). For the second, use
\(K_1(\widetilde\D_0)\subset c_0^{-1}PK_1(\D_0)\) and map an entropy cover through \(P\).
\end{proof}

\subsection{Dictionary residuals}

For \(r\in X^*\), introduce the dictionary dual seminorm
\[
  \norm r_{\D^*}=\sup_{g\in\D_X}\abs{\ip r g}.
\]
This quantity measures the derivative in the available dictionary directions. It is a seminorm;
if the dictionary has dense linear span, its vanishing implies that the derivative is zero.

\subsection{Ideal and finite-pool Chebyshev greedy algorithms}

The ideal CGA combines weak selection with exact minimization over the enlarged space.

\begin{algorithm}[htbp]
\caption{Continuous weak Chebyshev greedy algorithm}
\label{alg:continuous}
Assume \(\E(0)<\infty\), choose weakness numbers \(t_m\in(0,1]\), and set
\(G_0=0\), \(\Phi_0=\{0\}\). Given
\(G_m\) and \(\Phi_m=\spn\{g_1,\ldots,g_m\}\):
\begin{enumerate}[label=\textup{(\roman*)}]
  \item if \(\norm{\E'(G_m)}_{\D^*}=0\), stop; otherwise choose \(g_{m+1}\in\D_X\) with descent sign such that
  \[
  -\ip{\E'(G_m)}{g_{m+1}}
  \geq t_m\norm{\E'(G_m)}_{\D^*};
  \]
  \item set \(\Phi_{m+1}=\spn(\Phi_m,g_{m+1})\) and compute
  \[
    G_{m+1}=\argmin_{v\in\Phi_{m+1}}\E(v).
  \]
\end{enumerate}
\end{algorithm}

For a general continuous dictionary the supremum need not be attained. Thus \(t_m<1\) is understood as
approximate-supremum selection. The endpoint \(t_m=1\) is used only when attainment has been proved,
for example on a compact nondegenerate parameter set with continuous score. The finite-pool method
stops, with a recorded reason, if the unused pool is empty, every remaining score is below the
stationarity threshold, every candidate fails the rank or acceptance test, or the attempt budget is
reached. In an unconstrained active-space solve, changing an atom's sign leaves the span unchanged;
the descent sign required by the proof is therefore not an additional algorithmic variable.

Exact reoptimization yields the variational orthogonality
\begin{equation}
 \ip{\E'(G_m)}v=0\qquad(v\in\Phi_m).
 \label{eq:orthogonality}
\end{equation}
The second algorithm describes the computed object. It uses a finite pool \(\Ppool_m\), an evaluated
score \(\widetilde s_m\), and a correction computed to a finite tolerance.

\begin{algorithm}[htbp]
\caption{Finite-pool inexact CGA}
\label{alg:finite}
Set \(\widehat G_0=0\), \(\widehat\Phi_0=\{0\}\). At step \(m\), form a finite admissible pool
\(\Ppool_m\subset\D\). Let
\[
 s_m(g)=\abs{\ip{\E'(\widehat G_m)}g},
 \qquad \widetilde s_m(g)\geq0.
\]
Choose \(t_m^{\rm pool}\in(0,1]\) and \(\widehat g_{m+1}\in\Ppool_m\) satisfying
\[
 \widetilde s_m(\widehat g_{m+1})
 \geq t_m^{\rm pool}\max_{g\in\Ppool_m}\widetilde s_m(g),
\]
give it the true descent sign, and define
\(\widehat\Phi_{m+1}=\spn(\widehat\Phi_m,\widehat g_{m+1})\). Compute a coefficient vector whose state
obeys
\[
 \E(\widehat G_{m+1})
 \leq\inf_{v\in\widehat\Phi_{m+1}}\E(v)+\varepsilon_m^{\rm opt},
 \qquad \varepsilon_m^{\rm opt}\geq0.
\]
\end{algorithm}

\subsection{Finite-pool error quantities}
\label{sec:controlled-selection}

The continuous-dictionary supremum, the exact maximum over the finite pool, and the
numerically evaluated maximum are distinct quantities. We impose the following
score-error and pool-coverage conditions on any step for which a positive coverage
bound is used; finiteness of the pool alone does not imply a uniform positive lower bound.
\begin{align}
 \sup_{g\in\Ppool_m}\abs{\widetilde s_m(g)-s_m(g)}&\leq\xi_m,
 \label{eq:score-error}\\
 \max_{g\in\Ppool_m}s_m(g)&\geq
 \tau_m\sup_{g\in\D}s_m(g),\qquad0<\tau_m\leq1.
 \label{eq:pool-coverage}
\end{align}

\begin{lemma}[True score selected from an evaluated pool]
\label{lem:true-score}
Under \cref{eq:score-error,eq:pool-coverage},
\begin{equation}
 s_m(\widehat g_{m+1})\geq
 t_m^{\rm pool}\tau_m\norm{\E'(\widehat G_m)}_{\D^*}
 -(1+t_m^{\rm pool})\xi_m.
\label{eq:true-score-lower}
\end{equation}
\end{lemma}

\begin{proof}
Insert the score error before and after the finite-pool maximization:
\begin{align*}
s_m(\widehat g_{m+1})
&\geq\widetilde s_m(\widehat g_{m+1})-\xi_m\\
&\geq t_m^{\rm pool}\max_{g\in\Ppool_m}\widetilde s_m(g)-\xi_m\\
&\geq t_m^{\rm pool}\max_{g\in\Ppool_m}s_m(g)
 -(1+t_m^{\rm pool})\xi_m,
\end{align*}
and use \cref{eq:pool-coverage}.
\end{proof}

If the computed correction minimizes a discrete training objective $\E_h$, a single terminal comparison
between $\E_h$ and the analysis energy $\E$ is not a uniform error bound. The appropriate deterministic statement is
the following.

Let \(G_{m+1}^{\sharp}\) denote an exact minimizer of \(\E\) over
\(\widehat\Phi_{m+1}\). For an energy level \(\Lambda_m^E\) chosen so that both
\(G_{m+1}^{\sharp}\) and the computed state \(\widehat G_{m+1}\) are included, define the relevant sublevel
set and its active-space restriction by
\begin{equation}
 \mathcal L_m=\{v\in X:\E(v)\leq\Lambda_m^E\},
 \qquad \mathcal L_m^{\rm act}=\widehat\Phi_{m+1}\cap\mathcal L_m.
 \label{eq:active-sublevel}
\end{equation}
Thus \(\mathcal L_m^{\rm act}\) is the part of the current finite-dimensional trial space on which the
training and analysis energies must be compared; in particular,
\(\inf_{v\in\mathcal L_m^{\rm act}}\E(v)=\E(G_{m+1}^{\sharp})\).

\begin{lemma}[Transfer from the discrete energy to the analysis energy]
\label{lem:energy-transfer}
Suppose on the active-space sublevel set \(\mathcal L_m^{\rm act}\) in
\cref{eq:active-sublevel} that
\[
 \sup_{v\in\mathcal L_m^{\rm act}}
 \abs{\E_h(v)-\E(v)}\leq\zeta_m^E.
\]
If the computed state satisfies
\[
 \E_h(\widehat G_{m+1})\leq
 \inf_{v\in\mathcal L_m^{\rm act}}\E_h(v)+\varepsilon_m^{\rm opt,h},
\]
then \Cref{alg:finite} holds with
\(\varepsilon_m^{\rm opt}\leq\varepsilon_m^{\rm opt,h}+2\zeta_m^E\).
\end{lemma}

\begin{proof}
Apply the uniform energy difference once at \(\widehat G_{m+1}\) and once at
\(G_{m+1}^{\sharp}\). Both states belong to \(\mathcal L_m^{\rm act}\), and the latter minimizes the analysis
energy over the full active space.
\end{proof}
